\scanpage{126}
\mspage{113}
Mais le diagramme (6.3.6), où la flèche
en traits discontinus $h'_{I'}$ est déjà acquise,
et $h_0$ en pointillés reste à trouver, donne
naissance au diagramme, dans lequel
est \unsure{mis|mise} en évidence la flèche \uline{$\widetilde i\in\Phi$}
qui exprime l'hypothèse $i\in C(I)$,
compte tenu de $i'\in C(I')$ :

\diagram{\((6.3.9)\)\quad
\begin{tikzcd}[ampersand replacement=\&,column sep=small,row sep=small]
 A'_{I'}\arrow[rr,"{\alpha_{I'}}"]\arrow[dr,"{i'_{I'}}"]\arrow[dd,"{\lambda_A}"'] \& \& X'_{I'}\arrow[dr,"{p_{I'}}"]\arrow[dd,"{\lambda_X}"{pos=.23}] \& \\
 \& B'_{I'}\arrow[rr,"{\beta_{I'}}",crossing over]\arrow[ur,dashed,"{h'_{I'}}"]\arrow[dd,"s"'{pos=.77},crossing over] \& \& Y'_{I'}\arrow[dd,"{\lambda_Y}"]\\
 A_0\arrow[rr,"{\alpha_0}"{pos=.78}]\arrow[dr,"r"']\arrow[ddr,dashed,"{i_0}"']\arrow[ur,phantom,"{\text{\uline{cocart}}}"{pos=.35}] \& \& X_0\arrow[dr,"{p_0\in\Psi}"] \& \\
 \& \widetilde B\arrow[rr,preaction={draw,solid,line width=.35pt,-},dashed,very thick,"{\widetilde\beta_0}"]\arrow[ur,dashed,thick,"{\widetilde h_0}"]\arrow[d,"{\deleted{\unsure{i_0}}\widetilde i\in\Phi}"'] \& \& Y_0\\
 \& B_0\arrow[uur,dotted,"{h_0}"]\arrow[urr,"{\beta_0}"'] \& \&
\end{tikzcd}}

où la flèche $\widetilde h_0$ en traits discontinus,
de $\widetilde B\overset{\mathrm{def}}{=}A_0\amalg_{A'_{I'}}B'_{I'}$ dans $X_0$, est définie
par les conditions de commutativité
\[
 \widetilde h_0r=\alpha_0,\qquad \unsure{\bar h|\widetilde h}_0s=\lambda_X h'_{I'}\tag*{(6.3.10)}
\]
ce qui est licite, vu que
\[
 \begin{aligned}
 \alpha_0\lambda_A &=
 \underset{\Vert}{\underline{(\lambda_X h'_{I'})i'_{I'}}}\qquad\deleted{\unsure{=}}\\[-2pt]
 &\phantom{=}\underbrace{\lambda_X\bigl(\underset{\alpha_{I'}}{\underline{h'_{I'}i'_{I'}}}\bigr)}_{\substack{\Vert\\\lambda_X\alpha_{I'}}}
 \end{aligned}
\]
% archive sheet number: 125

